\documentclass[10pt,letterpaper]{article}
\usepackage[margin=1in,headheight=14pt]{geometry}
\usepackage[T1]{fontenc}
\usepackage{lmodern}
\usepackage{amsmath,amssymb,amsthm,mathtools}
\usepackage{microtype,booktabs,xcolor,fancyhdr}
\definecolor{ink}{HTML}{19354A}
\usepackage[colorlinks=true,linkcolor=ink,urlcolor=ink,citecolor=ink]{hyperref}
\hypersetup{pdftitle={Vertex color-coding for simple paths},pdfauthor={Akhil Ravuri}}
\pagestyle{fancy}\fancyhf{}
\fancyhead[L]{\small\textsc{Vertex color-coding}}
\fancyhead[R]{\small Akhil Ravuri}\fancyfoot[C]{\thepage}
\renewcommand{\headrulewidth}{0.3pt}
\setlength{\parskip}{0.2em}\setlength{\emergencystretch}{1em}
\title{\vspace{-1.5em}\LARGE\bfseries Vertex color-coding\\for simple paths}
\author{Akhil Ravuri}\date{}
\begin{document}
\maketitle\vspace{-1.5em}
\begin{abstract}
We describe a teaching implementation of the standard color-coding dynamic
program for colorful simple paths. A fixed-path counting argument gives the
one-trial success guarantee. Exhaustive checks on every labeled simple graph
of order at most four compare the implementation with ordered-vertex path
enumeration. This is an exposition and implementation reproduction of a known
method, without a novelty claim.
\end{abstract}

\section{Paths and color sets}
Let $G=(V,E)$ be a finite simple undirected graph with $|V|=n$, and let
$1\le k\le n$. Assign each \emph{vertex} a color in $\{0,\ldots,k-1\}$.
A colorful path uses distinct colors. A $k$-vertex colorful path consequently
uses every color and has $k-1$ edges. These vertex colors differ from the edge
colors in the repository's rainbow-matching experiments.

The method follows Alon, Yuster, and Zwick~\cite[Section 3 and Lemma 3.1]{AYZ}.
For a fixed coloring $c$, define $D[S,v]$ to mean that a colorful path ending
at $v$ uses exactly the color set $S$. Initialize
\[
 D[\{c(v)\},v]=\mathrm{true}\qquad(v\in V).
\]
For every reachable state and neighbor $w$ of $v$ with $c(w)\notin S$, set
\[
 D[S\cup\{c(w)\},w]=\mathrm{true}.
\]
The implementation represents $S$ by a bit mask and stores one explicit path
for each reachable state. Adding a missing bit strictly increases the mask,
so processing masks in increasing numerical order is valid.

\section{Why the dynamic program is correct}
Every initialized state is a one-vertex path. Every transition appends an
adjacent vertex with a new color. Such a vertex cannot already occur in the
stored path, since its color is absent. Inductively, each stored witness is a
simple colorful path with the claimed endpoint and color set.

Conversely, removing the last vertex of any colorful path leaves a shorter
colorful path whose state precedes it. Applying the corresponding extension
recovers its endpoint and color set. Storing only one witness for a state
does not lose future possibilities: any neighbor whose color is new is
absent from every witness with that color set, regardless of which vertices
the witnesses otherwise use. Thus equal endpoint and color-set states admit
the same extensions. A full-mask state exists exactly when a colorful
$k$-vertex path exists.

With explicit tuple copying, the teaching implementation takes
\[
 O\bigl(k2^k(n+|E|)\bigr)\text{ time},\qquad
 O(k2^k n)\text{ space}.
\]
There are at most $2^k n$ states, each holding at most $k$ vertices. Copying a
witness costs $O(k)$ when extending it. These bounds describe the current
implementation, rather than every optimized variant of color-coding.

\section{Success probability for an uncolored path}
Now color vertices independently and uniformly from $k$ colors. For a fixed
$k$-vertex path, exactly $k!$ of its $k^k$ color assignments are injective.
Its colorful probability is therefore
\[
 p_k=\frac{k!}{k^k}.
\]
If the uncolored graph contains such a path, the dynamic program finds some
path with probability at least $p_k$: making that one fixed path colorful
already suffices. Other paths can increase the success probability.
A failed trial does not certify the absence of an uncolored path. The
experiment below computes exact success fractions by enumerating assignments;
it does not estimate probabilities by random trials.

\section{Exhaustive finite verification}
For every $n=1,\ldots,4$, the experiment enumerates all
$2^{\binom n2}$ labeled simple undirected graphs. For each $k=1,\ldots,n$,
it tests every one of the $k^n$ assignments. The total is
\[
 \sum_{n=1}^{4}2^{\binom n2}\sum_{k=1}^{n}k^n=22{,}955.
\]
For each assignment, an independent oracle enumerates ordered $k$-tuples of
distinct vertices, checks adjacency, and tests colorfulness. It agrees with
the dynamic program in all 22,955 cases. Every returned witness is checked
for its vertex count, distinct vertices, adjacency, and distinct colors.

For each tested graph containing a $k$-vertex path, the number of successful
assignments divided by $k^n$ is checked to be at least $k!/k^k$ using exact
fractions. This tests the implementation on its stated finite domain; the
general correctness and fixed-path arguments are separate from the audit.
The code also accepts directed adjacency, but this exhaustive experiment
covers simple undirected graphs only.

\subsection*{Reproduction and scope}
From the repository root, with Python 3.11+ and its standard library, run
\begin{quote}\small\ttfamily
python -m experiments.color\_coding\_small\_cases\par
python -m experiments.reproduce --check results/data/small\_cases.json
\end{quote}
The routine is \path{src/color_coding.py}; the experiment and saved JSON
record exact assignment counts and success fractions. This is not an
implementation of the cited paper's full algorithm collection or its
perfect-hash derandomization. No new path algorithm or general empirical
performance claim is made. The
\href{https://github.com/Akhil-rav/discrete-math-lab/blob/main/notes/literature/source_records.md}{source records}
identify the author manuscript consulted.

\begin{thebibliography}{1}
\bibitem{AYZ} N. Alon, R. Yuster, and U. Zwick, \emph{Color-coding},
author manuscript, revised March 9, 1995.
\href{https://www.math.tau.ac.il/~nogaa/PDFS/col5.pdf}{Author-hosted manuscript}.
Section 3 and Lemma 3.1 supply the colorful-path method and fixed-path
probability argument used here.
\end{thebibliography}
\end{document}
